\chapter{简单命令}
\label{cha:intro}
这一章介绍一些基本的 \LaTeX 命令，包括字体、数学公式、表格、插图、参考文献等。

\section{封面}
封面的例子请参看 \verb|cover.tex| 。里面的命令非常简单，一看即会。
姓名、论文题目、导师姓名职称等信息在 \verb|name.tex| 中，自己填一下。

\section{字体}
\begin{center}
\begin{tabular}{lc}
\verb|\song| & \song 宋体\\
\verb|\kai|  & \kai  楷书\\
\verb|\hei|  & \hei  黑体\\
\verb|\fs|   & \fs   仿宋\\
\verb|\li|   & \li   隶书\\
\verb|\yy|   & \yy   幼圆\\
\verb|\zs|   & \zs   中宋\\
\end{tabular}
\end{center}

封面字体宋体加粗效果不好看，采用华文中宋字体。
字体下载请点击下面的链接或自行搜索该字体：\\
\url{https://jianguoyun.com/c/pubfile/Lye3gYc3po3MwHB-Cd2j9Vj-3EH6MG1XbiAmcxZzvak}
\section{数学}
自己找本书看吧。推荐 lshort 和 lnotes 这两本。\\
\url{http://www.dralpha.com/chs/tech/lnotes2.pdf}\\
\url{http://ftp.ctex.org/mirrors/CTAN/info/lshort/chinese/lshort-zh-cn.pdf}

\section{表格}
自己找本书看吧。

\section{插图}
插图命令：
\begin{verbatim}
\begin{figure}[htbp]
\centerline{\includegraphics[width=1.97cm]{logo}}
\caption{南京大学校徽}
\label{fig:logo}
\end{figure}
\end{verbatim}

\begin{figure}[htbp]
\centerline{\includegraphics[width=1.97cm]{logo}}
\caption{南京大学校徽}
\label{fig:logo}
\end{figure}

\begin{figure}[htbp]
\centerline{\includegraphics[height=2cm]{name}}
\caption{南京大学毛体校名}
\label{fig:name}
\end{figure}


\section{参考文献}
Casting Metal Nanowires Within Discrete Self-Assembled Peptide Nanotubes{\cite{casting}}

Blue Luminescence Based on Quantum Confinement at Peptide Nanotubes{\cite{qc}}

\section{编译方法}
采用 \texttt{xelatex} 编译，流程：
\begin{verbatim}
xelatex main
bibtex  main
xelatex main
xelatex main
\end{verbatim}
对于 Linux 用户，我写了个简单的 \texttt{Makefile}，在当前目录终端下输入 \texttt{make} 即可。

